%% =====================================================================
%%  A rotational sub-cell descriptor for cross-topology stiffness prediction
%%
%%  COMPLETE REWRITE. The earlier manuscript ("The target, not the learner")
%%  was computed on a dataset later found to be 61.6% duplicated. Every number
%%  here is recomputed on MetaXFam-D, a verified duplicate-free benchmark, and
%%  the central claim has changed as a result. See Sec. 6.
%%
%%  BUILD: pdflatex x3
%%  Template: stock article class. No journal class is emulated.
%% =====================================================================
\documentclass[11pt,a4paper]{article}
\usepackage[T1]{fontenc}
\usepackage[utf8]{inputenc}
\usepackage{mathptmx}
\usepackage{courier}
\usepackage[margin=2.4cm]{geometry}
\usepackage{amsmath,amssymb}
\usepackage{graphicx}
\usepackage{booktabs}
\usepackage{caption}
\usepackage{authblk}
\usepackage[hidelinks]{hyperref}
\usepackage[numbers,sort&compress]{natbib}
\usepackage{microtype}
\sloppy
\captionsetup{font=small,labelfont=bf}
\newcommand{\CH}{\mathbf{C}^{H}}
\newcommand{\Rtwo}{R^{2}}
\renewcommand{\thetable}{\Roman{table}}

\title{\bfseries A rotational sub-cell descriptor gives stable cross-topology\\
prediction of effective stiffness where pixel surrogates do not}

\author[1]{David Mashiah\thanks{davidmashiah@mail.tau.ac.il \quad ORCID 0009-0004-4684-955X}}
\affil[1]{School of Mechanical Engineering, Tel Aviv University, Tel Aviv 6997801, Israel}
\date{\today}

\begin{document}
\maketitle

\begin{abstract}
\noindent
Machine-learned surrogates for the homogenized stiffness of periodic metamaterials are
usually evaluated by splitting a pooled cell dataset at random. On a benchmark containing
several unit-cell topology families this measures interpolation within known families, not
the transfer a practitioner needs. We report three results on a purpose-built,
duplicate-free 18-family benchmark. First, the protocol matters more than the model: the
same surrogate scores $\Rtwo$ between $0.90$ and $0.96$ under a random split and between
$-0.73$ and $+0.14$ when whole topology families are held out, a gap of $0.58$ to $1.63$
across four stiffness components. Second, a three-number descriptor of how the softest
sub-cell elastodynamic mode divides its motion between rotation, dilatation and shear --
obtained from the same periodic operator that produces the stiffness tensor, at a quarter
of the cost of one homogenization -- is the only representation we tested that never fails:
across twelve component-by-learner combinations its held-out $\Rtwo$ stays within
$[+0.02, +0.41]$, whereas $2304$ raw pixels span $[-42.99, +0.34]$ and a published
geometric spectral descriptor spans $[-1.56,+0.33]$. Third, its advantage is not uniform
and we report where it fails: it beats the geometric spectrum and the free solid-fraction
control on the shear-coupled components $C_{12}$ and $C_{22}$ for every learner tested, but
does not beat solid fraction at all on $C_{11}$. That pattern is what a
rotational-mode explanation predicts, and it is the reason we do not claim a general
advantage. We also document a defect in the dataset the earlier version of this work used:
rasterizing low-parameter shape families onto a fixed grid produces massive duplication,
$61.6\%$ in that case, which inflates in-distribution scores and depresses pixel baselines.
\end{abstract}

\noindent\textbf{Keywords:} computational homogenization; metamaterials; machine-learning
surrogates; out-of-distribution generalization; micropolar effects; benchmark design

\section{Introduction}

Numerical homogenization returns an effective stiffness $\CH$ for a periodic microstructure
at the cost of one cell problem per cell. Because that cost is prohibitive inside a design
loop, a large literature replaces it with a learned surrogate mapping geometry to $\CH$,
with reported accuracy of $\Rtwo \ge 0.98$ \citep{bastek2022,kumar2020}.

That accuracy is measured within a design space. Whether it survives a change of
unit-cell topology is a separate question, and a recent broad study of materials-property
models concludes that most out-of-distribution evaluations in the field measure
interpolation rather than extrapolation \citep{li2025}. Work on generalizable
microstructure--property maps across microstructure classes reaches a related conclusion
and additionally reports that a distributional distance predicts when transfer will succeed
\citep{liu2026}.

This paper asks a narrower question with a controlled benchmark: when a surrogate is
required to predict $\CH$ for a topology family it has never seen, how much does the
evaluation protocol matter, and is there a cheap physically-motivated representation that
degrades gracefully rather than catastrophically?

Our starting point is that $\CH$ is obtained by volume-averaging the microscopic stress.
Averaging is a projection: it retains the two rigid-translation modes of the cell and
discards every other sub-cell degree of freedom. Roberts \citep{roberts2025} constructs
multi-continuum micromorphic homogenizations by retaining precisely those discarded modes,
and the relaxed-micromorphic literature makes the same point constitutively --- a
first-order Cauchy continuum has no slot for the cell's rotational degree of freedom
\citep{neff2014}. We use those modes not as a modelling target but as an inexpensive
descriptor, and we test where that descriptor helps and where it does not.

\paragraph{Contributions.}
(i) MetaXFam-D, an 18-family, 3096-cell benchmark verified to contain no duplicate
geometries, with family-disjoint evaluation (\S\ref{sec:bench}).
(ii) A measurement of how much a random split overstates cross-topology performance
(\S\ref{sec:gap}).
(iii) A three-number rotational sub-cell descriptor that is the only representation tested
whose held-out $\Rtwo$ never becomes negative (\S\ref{sec:stab}).
(iv) An explicit map of where it beats the baselines and where it does not, with confidence
intervals and a pre-registered aggregation rule (\S\ref{sec:where}).
(v) A dataset-construction caution: rasterization collapse (\S\ref{sec:dup}).

\paragraph{What we do not claim.} We do not claim the descriptor generally outperforms a
pixel surrogate; on a duplicate-free benchmark it does not (\S\ref{sec:where}). We do not
claim non-injectivity of homogenization as novel; it is established \citep{kumar2021}.
We do not claim that distributional criteria fail to predict transfer --- on this benchmark
they partly succeed, consistent with \citep{liu2026}.

\section{Benchmark}
\label{sec:bench}

Cells are $48\times48$ binary grids, plane strain, solid $E=1$, $\nu=0.3$, void
$E=10^{-6}$, with periodicity imposed by identifying opposite-edge nodes. $\CH$ follows
from volume-averaging element stresses under three unit macroscopic strains
\citep{andreassen2014}; the solver passes six validation tests including the exact rank-1
laminate and Voigt--Reuss bracketing.

Every cell satisfies three conditions: solid fraction in $[0.45,0.75]$; the solid phase
percolates in both directions; and its rasterization is unique within its family. The
benchmark comprises 18 families $\times$ 172 cells $=3096$ cells, verified to contain
$3096$ distinct geometries and no cross-family duplicates. Family median solid fractions
span $0.579$--$0.695$, so density is controlled and a cross-family error is a topology
effect rather than a density extrapolation.

\subsection{Rasterization collapse, and why the benchmark had to be rebuilt}
\label{sec:dup}

An earlier version of this dataset sampled continuous shape parameters and kept 600 cells
per family. Auditing it revealed that its $13{,}200$ cells comprise only $5072$ distinct
rasterizations --- a duplicate rate of $61.6\%$. Low-parameter families are the worst
affected: one lattice family admits only \emph{three} distinct $48\times48$ rasterizations
inside the density band, and rejects essentially every further draw as already seen.

The mechanism is confirmed by resolution dependence (Fig.~\ref{fig:dup}): duplication rises
from $72.8\%$ at $48\times48$ to $98.4\%$ at $12\times12$. Coarser rasterization admits
fewer distinct images. The consequence is severe: under a random $75/25$ split of that
dataset, $67.7\%$ of test cells have their exact geometry in the training set.

This is not a general property of pixelized metamaterial datasets. The published 2D subsets
of METASET \citep{chan2021}, audited with the same script, show $0\%$ duplication at
$50\times50$; that dataset selects for diversity explicitly and uses richer shape
descriptions. The defect belongs to low-parameter parametric generators. We report it
because the check is three lines of code and, in our case, it changed a headline
conclusion.

\begin{figure}[t]
\centering
\includegraphics[width=0.42\textwidth]{f4_dup.png}
\caption{Duplication against rasterization resolution, measured on an eight-family probe.
Distinct shape parameters collapse onto identical pixel images as resolution falls.}
\label{fig:dup}
\end{figure}

\subsection{Evaluation protocol}

Training subsets of $K=4$ families are drawn at random; 30 subsets $\times$ 2 model seeds
gives 60 paired evaluations per arm, scored as pooled $\Rtwo$ over the 14 held-out
families. Three surrogate classes are used throughout --- random forest, ridge regression
with cross-validated regularization, and a two-layer perceptron --- because
representation comparisons that hold for one learner need not hold for another.

\paragraph{Aggregation rule, fixed in advance.} A claim is reported as supported only if it
holds for \emph{every} learner on a given component, with a bootstrap $95\%$ confidence
interval on the paired median gain excluding zero. We state this because an earlier
analysis of ours summarized on a single learner and reached a different headline from the
same measurements.

\section{The protocol gap}
\label{sec:gap}

Table~\ref{tab:gap} compares three protocols at matched training-set size: a random split
within the training families, a random split pooled over all families, and a
family-disjoint split. The same model and the same number of training cells are used
throughout; only the split changes.

\begin{table}[t]
\centering\small
\caption{Held-out $\Rtwo$ by protocol, ridge regression on raw pixels, median over 20
training subsets. Only the split protocol differs.}
\label{tab:gap}
\begin{tabular}{lrrrr}
\toprule
component & random, same families & random, pooled & family-disjoint & gap \\
\midrule
$C_{11}$ & $+0.961$ & $+0.859$ & $+0.142$ & $0.819$ \\
$C_{12}$ & $+0.897$ & $+0.745$ & $-0.728$ & $1.625$ \\
$C_{22}$ & $+0.947$ & $+0.829$ & $-0.234$ & $1.181$ \\
$C_{33}$ & $+0.912$ & $+0.752$ & $-0.416$ & $1.328$ \\
\bottomrule
\end{tabular}
\end{table}

The pooled random split --- the standard practice when a dataset contains several families
--- still reports $\Rtwo$ between $0.75$ and $0.86$ while the same surrogate is at or below
zero on unseen topologies (Fig.~\ref{fig:gap}). A practitioner following that protocol
would report a strong result and ship a model that does not transfer.

\begin{figure}[t]
\centering
\includegraphics[width=0.55\textwidth]{f3_splitgap.png}
\caption{The same surrogate under three evaluation protocols on the clean benchmark.}
\label{fig:gap}
\end{figure}

\section{The descriptor}

On the same periodic operator, with a consistent mass matrix, we solve
$\mathbf{K}\boldsymbol{\phi}=\omega^{2}\mathbf{M}\boldsymbol{\phi}$. Periodicity admits
exactly two rigid modes; the modes above them are what stress-averaging discards. For each
eigenvector the displacement gradient splits pointwise and exactly into dilatation,
deviatoric shear and rotation,
\begin{equation}
|\nabla u|^{2}=\tfrac12 d^{2}
+\tfrac12\!\left[(\varepsilon_{xx}-\varepsilon_{yy})^{2}+\gamma_{xy}^{2}\right]+2w^{2},
\qquad w=\tfrac12(\partial_x v-\partial_y u),
\end{equation}
and each term is integrated over the cell weighted by element stiffness and normalized to
sum to one. We call the result the \emph{mode character} of that eigenvector, and use the
character of the softest non-rigid mode: three numbers.

An eigenvector inside a degenerate cluster is defined only up to a rotation within the
cluster and the character is quadratic in the eigenvector, so per-eigenvector character is
basis-dependent. Every four-fold-symmetric family has a degenerate second mode. We
therefore average over the cluster, which is basis-invariant, and use a deterministic
eigensolver start.

\paragraph{Validation.} On a homogeneous cell the exact periodic modes are plane waves: the
lowest non-rigid cluster is the shear wave at $\omega=2\pi\sqrt{\mu}=3.8967$ with character
exactly $(0,\tfrac12,\tfrac12)$; the pressure wave has $(\tfrac12,\tfrac12,0)$. The
extractor reproduces the fundamental to $7.1\times10^{-4}$, converges as $h^{2}$, reproduces
both exact characters to $10^{-6}$, and is invariant to a random rotation within a
degenerate cluster to $10^{-14}$. Ten tests pass.

\paragraph{Cost.} Computed on a $24\times24$ coarsened cell, the eigensolve costs $0.25\times$
one $48\times48$ homogenization.

\section{Stability across surrogate classes}
\label{sec:stab}

The clearest result concerns not average accuracy but the absence of failure.
Table~\ref{tab:stab} reports, for each representation, the range of median held-out $\Rtwo$
over all twelve component-by-learner combinations.

\begin{table}[t]
\centering\small
\caption{Range of median held-out $\Rtwo$ across four components $\times$ three learners.}
\label{tab:stab}
\begin{tabular}{lrrr}
\toprule
representation & dim & min & max \\
\midrule
raw pixels & 2304 & $-42.99$ & $+0.34$ \\
solid fraction $+$ shape-DNA & 7 & $-1.56$ & $+0.33$ \\
solid fraction & 1 & $+0.04$ & $+0.24$ \\
\textbf{solid fraction $+$ mode-2 character} & \textbf{4} & $\mathbf{+0.02}$ & $\mathbf{+0.41}$ \\
\bottomrule
\end{tabular}
\end{table}

The pixel surrogate attains the single best cell in the table and also the worst by a factor
of a hundred: under a perceptron it reaches $-42.99$ on $C_{33}$ and $-18.12$ on $C_{12}$.
The four-number representation never becomes negative in any of the twelve cells
(Fig.~\ref{fig:stab}). Solid fraction alone is similarly stable but reaches only $+0.24$.

We regard this as the practically useful statement. A practitioner choosing a surrogate for
an unseen topology cannot know in advance which learner will be well matched to the data;
a representation that is mediocre-but-never-catastrophic is more valuable than one whose
expected performance is higher but whose variance across model classes spans forty units of
$\Rtwo$.

\begin{figure}[t]
\centering
\includegraphics[width=0.62\textwidth]{f1_stability.png}
\caption{Each point is one component $\times$ learner combination.}
\label{fig:stab}
\end{figure}

\section{Where the descriptor helps, and where it does not}
\label{sec:where}

Applying the pre-registered rule, none of the three pairwise claims holds across all four
components (Fig.~\ref{fig:gains}). The pattern is systematic.

\paragraph{Against the free control (solid fraction).} The descriptor wins on $C_{22}$ for
every learner ($+0.315$, $+0.212$, $+0.195$; all intervals clear of zero; $90$--$95\%$ of
paired evaluations) and on $C_{12}$ for two of three. On $C_{11}$ it wins for none: the
gains are $+0.046$, $-0.037$, $-0.060$, with intervals containing zero.

\paragraph{Against the published geometric spectrum.} Shape-DNA
\citep{wang2020} is the Laplace--Beltrami spectrum of the cell --- a spectrum, rotation
invariant, purely geometric, and the closest prior art. The elastic descriptor beats it in
ten of twelve cells, decisively on $C_{22}$ ($+0.161$ to $+0.656$, $90$--$100\%$ of
evaluations). It loses on $C_{11}$ under a random forest ($-0.164$, interval entirely
negative). So it is not the case that any spectrum transfers: the elastic one does on the
shear-coupled components and the geometric one largely does not.

\paragraph{Against raw pixels.} The descriptor wins under ridge and the perceptron on every
component, but under a random forest it is indistinguishable on $C_{22}$ and $C_{33}$ and
loses on $C_{11}$ ($-0.187$). An earlier version of this work, computed on the duplicated
dataset, reported a uniform $+0.364$ advantage; that advantage was an artifact of
duplication depressing the pixel baseline.

\begin{figure}[t]
\centering
\includegraphics[width=\textwidth]{f2_gains.png}
\caption{Paired gain of the four-number descriptor over each baseline, with bootstrap
$95\%$ intervals. Blue: interval excludes zero. Red: it does not. Values clipped for display.}
\label{fig:gains}
\end{figure}

\paragraph{Reading the pattern.} The advantage is concentrated in $C_{12}$ and $C_{22}$ and
absent in $C_{11}$. Under the rotational-mode account this is the expected signature rather
than an inconvenience: the discarded degree of freedom is a rotation, and rotation couples
to shear, so the components it should inform are the shear-coupled ones. A descriptor that
improved every component equally would in fact be harder to explain.

\section{Discussion}
\label{sec:disc}

\paragraph{What changed from the earlier version, and why.} A previous version of this work
reported that a low-dimensional physical feature broadly outperformed a pixel surrogate,
on a dataset later found to be $61.6\%$ duplicated. Recomputing every number on
duplicate-free data removed that conclusion: the pixel baseline recovers from $-0.008$ to
$+0.37$ on $C_{22}$ once duplicates are gone. We report this because the direction of the
bias is instructive --- duplication does not merely inflate in-distribution scores, it
\emph{depresses} high-dimensional baselines on cross-family evaluation, and so
systematically flatters low-dimensional competitors.

\paragraph{Relation to prior work.} Francis et al.\ \citep{francis2025} learn homogenized
micropolar constants with a deep material network, extrapolating across constitutive
behaviour and boundary conditions rather than topology. Hendriks et al.\
\citep{hendriks2025} address cross-topology generalization by making the architecture
equivariant to geometric similarity; that is complementary to changing what is measured, and
we do not compare against it here, which is the principal gap in this work. Liu et al.\
\citep{liu2026} report that a distributional distance predicts generalizability across
microstructure classes; on this benchmark we reproduce that qualitatively.

\paragraph{Limitations.} Two dimensions, linear elasticity, one contrast ratio, one void
phase. The benchmark is 3096 cells; six families from the original pool could not supply
enough distinct rasterizations and were re-parameterized, which changes what those families
denote. No comparison against an equivariant graph network. The descriptor does not make
cross-topology prediction accurate --- the best held-out $\Rtwo$ we report is $+0.41$.

\section*{Data and code availability}

The benchmark, generators, validated solver and duplication-audit tooling are archived at
Zenodo, \url{https://doi.org/10.5281/zenodo.22167146}, and at
\url{https://github.com/davidmashiah/metaxfam22}. The superseded duplicated dataset is
retained at \url{https://doi.org/10.5281/zenodo.21734948} because published work refers to
it. The sub-cell extractor and its validation suite are at
\url{https://doi.org/10.5281/zenodo.22137052} and
\url{https://github.com/davidmashiah/metaxfam-subcell}.

\section*{Declaration of generative AI in the writing process}

During the preparation of this work the author used a large language model as a coding and
drafting assistant: to implement and run analysis scripts and to draft and revise prose.
All numerical results were produced by the deposited code and checked by the author against
the raw result files. All scientific claims, their framing, and the decision to report each
negative result are the author's own. The author takes full responsibility for the content
of this publication.

\begin{thebibliography}{99}\small\raggedright

\bibitem{bastek2022} J.-H. Bastek, S. Kumar, B. Telgen, R. N. Glaesener, D. M. Kochmann,
Inverting the structure--property map of truss metamaterials by deep learning,
\emph{Proc.\ Natl.\ Acad.\ Sci.\ U.S.A.} \textbf{119}(1) e2111505119 (2022).
\url{https://doi.org/10.1073/pnas.2111505119}

\bibitem{kumar2020} S. Kumar, S. Tan, L. Zheng, D. M. Kochmann, Inverse-designed spinodoid
metamaterials, \emph{npj Comput.\ Mater.} \textbf{6} 73 (2020).
\url{https://doi.org/10.1038/s41524-020-0341-6}

\bibitem{li2025} K. Li, A. N. Rubungo, X. Lei et al., Probing out-of-distribution
generalization in machine learning for materials, \emph{Commun.\ Mater.} \textbf{6} 9
(2025). \url{https://doi.org/10.1038/s43246-024-00731-w}

\bibitem{liu2026} H. Liu, N. Baishnab, B. S. S. Pokuri, B. Ganapathysubramanian, O. Wodo,
Constructing generalizable microstructure--property maps across diverse microstructure
classes, \emph{MRS Communications} \textbf{16} 449--458 (2026).
\url{https://doi.org/10.1557/s43579-025-00873-z}

\bibitem{roberts2025} A. J. Roberts, Accurate families of multi-continuum micromorphic
homogenisations in multi-D space-time via dynamical systems theory, \emph{Trans.\ Math.\
Appl.} \textbf{9}(1) tnaf001 (2025). \url{https://doi.org/10.1093/imatrm/tnaf001}

\bibitem{neff2014} P. Neff, I.-D. Ghiba, A. Madeo, L. Placidi, G. Rosi, A unifying
perspective: the relaxed linear micromorphic continuum, \emph{Continuum Mech.\ Thermodyn.}
\textbf{26}(5) 639--681 (2014). \url{https://doi.org/10.1007/s00161-013-0322-9}

\bibitem{kumar2021} S. Kumar, D. M. Kochmann, What machine learning can do for
computational solid mechanics, arXiv:2109.08419 (2021); in \emph{Current Trends and Open
Problems in Computational Mechanics}, Springer (2022).

\bibitem{wang2020} L. Wang, Y.-C. Chan, Z. Liu, P. Zhu, W. Chen, Data-driven metamaterial
design with Laplace--Beltrami spectrum as ``shape-DNA'', \emph{Struct.\ Multidisc.\ Optim.}
\textbf{61} 2613--2628 (2020). \url{https://doi.org/10.1007/s00158-020-02523-5}

\bibitem{chan2021} Y.-C. Chan, F. Ahmed, L. Wang, W. Chen, METASET: exploring shape and
property spaces for data-driven metamaterials design, \emph{J. Mech.\ Des.} \textbf{143}(3)
031707 (2021).

\bibitem{francis2025} N. M. Francis, D. Shin, R. A. Lebensohn, F. Pourahmadian,
R. Dingreville, Micropolar deep material network, \emph{Comput.\ Methods Appl.\ Mech.\
Engrg.} \textbf{446} 118329 (2025). \url{https://doi.org/10.1016/j.cma.2025.118329}

\bibitem{hendriks2025} R. Hendriks, V. Menkovski, M. Do\v{s}k\'a\v{r}, M. G. D. Geers,
O. Roko\v{s}, Similarity-equivariant graph neural networks for homogenization of
metamaterials, \emph{Comput.\ Methods Appl.\ Mech.\ Engrg.} \textbf{439} 117867 (2025).
\url{https://doi.org/10.1016/j.cma.2025.117867}

\bibitem{andreassen2014} E. Andreassen, C. S. Andreasen, How to determine composite material
properties using numerical homogenization, \emph{Comput.\ Mater.\ Sci.} \textbf{83}
488--495 (2014).

\end{thebibliography}
\end{document}
